\section{Introduction}\label{sec:intro}

Language-model agents now act through tools. A support agent looks up a customer, issues a refund, or cancels an order; a coding agent reads files, runs commands and calls services. The Model Context Protocol (MCP)~\cite{mcpspec} has made this tool layer a common interface: one agent client can reach many servers, and one server can serve many clients. Agents are also moving from reading data to changing it, which raises the cost of a wrong action.

Most evaluation of agents measures whether a task is completed when the environment behaves~\cite{liu2023agentbench,yao2024taubench,zhou2023webarena,jimenez2023swebench}. Production tools do not always behave. They hang, return errors, return half a record, or return bytes that do not parse. What matters then is not only whether the task is completed, which may be impossible, but what the agent does instead. It may stop and say that the lookup failed, which is safe. It may retry the same call many times in quick succession, invent a plausible value to fill the gap, tell the user that a refund was issued when the refund call failed, or reach for a different tool, such as cancelling an order because refunding it did not work. These are behavioural failures. A task-success score does not separate them from honest failure, and a user is often the first to find them.

Distributed systems faced a similar gap and closed it with chaos engineering: deliberately injecting faults into a running system and checking that a stated steady-state behaviour holds~\cite{basiri2016chaos,principlesofchaos}. Recent work has begun to apply fault injection, and in places the chaos-engineering label, to agents (\S\ref{sec:background}). We argue that the discipline should become routine for tool-using agents, adapted to two facts about them. First, the component under test is a stochastic policy, so a single run is weak evidence and verdicts must be statistical. Second, the interesting outcome is a behaviour (what the agent called, how often, and what it told the user), so the steady-state hypothesis must be stated as checkable expectations on behaviour rather than as a service-level metric.

\paragraph{Approach.} We place a transparent fault-injecting proxy on the tool path, between an unmodified agent and its MCP servers. An experiment names a tool, a fault, the point at which the fault fires, and the expected safe behaviour. The proxy injects the fault and records every call; a grader then applies deterministic predicates to the recorded calls and to the agent's final answer. No model is used as a judge, so the same trace and answer always produce the same verdict. Each experiment is repeated, judged against a pass threshold, and paired with a no-fault control that runs the same task and the same checks.

\paragraph{Contributions.} This is a vision paper with a pilot. We contribute:
\begin{itemize}
\item a failure model for tool-using agents with two axes, injected faults and unsafe behaviours (\S\ref{sec:failure-model});
\item a method, in the tradition of chaos engineering and of recent agent fault-injection work, that checks declared behavioural expectations as deterministic predicates over the call trace and final answer (no model judge), with per-experiment pass thresholds over repeated runs and paired no-fault controls, so that a result can gate a change (\S\ref{sec:method});
\item AgentChaos, an open-source implementation built on switchyard, a hook-based MCP proxy core designed to be shared with other runtime tools (\S\ref{sec:implementation});
\item a pilot study of one agent with two models on four fault experiments and their four controls, 160 sessions in all (\S\ref{sec:pilot}); and
\item a research agenda, including the protocol for a full study (\S\ref{sec:agenda}).
\end{itemize}
